\section{Signed quadratic products and finite oscillations}
Put $D_R=4\tau-\tr C_{2\tau}^2$,
$D_J=4\tau-\tr\mathcal K_\tau^2$ and
$S=\tr(\mathcal K_\tau-\mathcal K_\tau C_{2\tau})$.
The operators $C_\pm=(C_{2\tau}\pm\mathcal K_\tau)/2$ are positive
concentrations, not the positive/negative spectral parts of $\mathcal K_\tau$.

\begin{proposition}[Signed quadratic constants]\label{prop:defect-identities}
\label{thm:cross-exact}\label{cor:signed-constants}
Exactly,
\begin{equation}\label{short:eq:signed-defect}
\tr(C_\pm-C_\pm^2)=\tfrac14(D_R+D_J\pm2S),\qquad
\tr(C_+C_-)=\tfrac14(D_J-D_R).
\end{equation}
For $P=UU^*$, $A=(1-P)R_\tau P$ and $B=(1-P)J_\tau P$,
$D_R=\|A\|_2^2$, $D_J=\|B\|_2^2$ and
$S=\Re\langle A,B\rangle_{\mathrm{HS}}$.
For every $\tau>0$,
\begin{equation}\label{eq:cross-exact}
S=\frac{4\log2}{\pi^2}\cos(2\pi\tau)+R_2,\qquad
|R_2|\le\frac3{\pi^3\tau}.
\end{equation}
Writing $c_+=\cos(\pi\tau)$, $c_-=\sin(\pi\tau)$, it follows that
\begin{align}
\tr(C_\pm-C_\pm^2)&=
\frac{3\log\tau+\log(2\pi^3)+3(1+\gamma)+16(\log2)c_\pm^4}{4\pi^2}
+r_\pm,
\label{eq:defect-plus}\\
\tr(C_+C_-)&=
\frac{\log(2\pi\tau)+1+\gamma+2(\log2)\cos(4\pi\tau)}{4\pi^2}
+r_0,
\label{eq:cross-product}
\end{align}
where $|r_\pm|\le21/(8\pi^3\tau)$ and $|r_0|\le9/(8\pi^3\tau)$.
\end{proposition}
\begin{proof}
Expand $C_\pm^2$ and use trace cyclicity. Since
$\Pi_\pm=(R_\tau\pm J_\tau)/2$ are projections, their compressed
defects are $\|(1-P)\Pi_\pm P\|_2^2=\|A\pm B\|_2^2/4$,
which also proves the Hilbert--Schmidt identities.
For \eqref{eq:cross-exact}, put $t=2\pi\tau$ and use $u=x-y,v=x+y$
in the product-kernel trace. Integrating first in $v$ gives
\[
\tr(\mathcal K_\tau C_{2\tau})=\frac4{\pi^2}
\left[\sin t\int_0^t\frac{\sin^2(2w)}{2w^2}dw
-\cos t\int_0^t\frac{\sin^2w\sin2w}{w^2}dw\right].
\]
The infinite integrals are $\pi/2$ and $\log2$
(Lemma~\ref{lem:two-integrals}); their tails have modulus at most
$1/(2t)$ and $1/t$. Subtracting from $\tr\mathcal K_\tau=2\sin t/\pi$
proves \eqref{eq:cross-exact}.
Insert the exact quadratic and prolate defect formulas in
\eqref{short:eq:signed-defect}.
With $c=\cos(2\pi\tau)$, the oscillating numerator is
$2(\log2)(2c^2-1)\pm8(\log2)c
=4(\log2)(1\pm c)^2-6\log2$.
The error sums are $(1/2+4+6)/4=21/8$ and $(1/2+4)/4=9/8$
in units $(\pi^3\tau)^{-1}$.
In particular
$D_R=\tr(C_+-C_+^2)+\tr(C_--C_-^2)-2\tr(C_+C_-)$.
\end{proof}

\begin{theorem}[Finite quadratic traces]\label{thm:hs-discrete}
Let $s_k=\sin^2(\pi\varepsilon k)/(\pi^2k^2)$, $s_0=\varepsilon^2$.
For $N\ge1$, $d=2N+1$, $0<\varepsilon<1/2$, exactly
\begin{equation}\label{eq:hs-finite-identity}
\tr A_N^2=\sum_{|k|\le2N}s_k\left[
2(d-|k|)+\frac{2\sin(2\pi\varepsilon(d-|k|))}
{\sin(2\pi\varepsilon)}\right].
\end{equation}
For each $\delta\in(0,1/2)$, uniformly in
$0<\varepsilon\le1/2-\delta$ and $N\ge1$,
\begin{equation}\label{eq:hs-discrete}
2d\varepsilon-\tr A_N^2=
\frac2{\pi^2}\left[\log(4N\sin(\pi\varepsilon))+1+\gamma
 +(\log2)\cos(2\pi\varepsilon d)\right]
 +O_\delta\left(\varepsilon^2+\frac1{N\varepsilon}\right).
\end{equation}
For $\Pi_W=C_N(2W)$, uniformly in $N\ge1$, $0<W\le1/4$,
\begin{equation}\label{eq:dpss-discrete}
\tr(\Pi_W-\Pi_W^2)=
\frac{\log(4N\sin(2\pi W))+1+\gamma}{\pi^2}
 +O((NW)^{-1}),
\end{equation}
with an absolute constant, including the endpoint $W=1/4$.
\end{theorem}
\begin{proof}
Put $K=2N$, $\theta=2\pi\varepsilon$, $h=\pi\varepsilon$.
The entrywise square is $4s_{m-n}\cos^2(h(m+n))$.
For fixed $k=m-n$, the $M=d-|k|$ sums $m+n$ form the progression
$-(M-1),-(M-3),\ldots,M-1$; its cosine sum is
$\sin(M\theta)/\sin\theta$, proving \eqref{eq:hs-finite-identity}.
Write its unmodulated and modulated parts as $\Sigma_1,\Sigma_2$.
Parseval for an arc gives $\sum_{k\in\mathbb Z}s_k=\varepsilon$, so
\[
2d\varepsilon-\Sigma_1=4d\sum_{k>K}s_k+4\sum_{k=1}^Kks_k.
\]
For decreasing $b_k\downarrow0$, geometric and Abel summation give
$|\sum_{k>K}b_k\cos(k\theta)|\le b_{K+1}/\sin(\theta/2)$.
Using $b_k=k^{-2},k^{-1}$, $\sin(\pi\varepsilon)\ge2\varepsilon$,
$\sum_{k>K}k^{-2}=K^{-1}+O(K^{-2})$ and
$H_K=\log K+\gamma+O(K^{-1})$, one obtains
\[
4d\sum_{k>K}s_k=\frac2{\pi^2}+O((N\varepsilon)^{-1}),\quad
4\sum_{k=1}^Kks_k=
\frac2{\pi^2}[\log(4N\sin(\pi\varepsilon))+\gamma]
 +O((N\varepsilon)^{-1}).
\]
Here $\sum_{k\ge1}\cos(k\theta)/k=-\log(2\sin(\theta/2))$ follows
by Abel limiting the power series for $-\log(1-re^{i\theta})$;
the displayed tail estimate justifies that limit.

Set $g_c(w)=\sin^2w\cos(2w)/w^2$,
$g_s(w)=\sin^2w\sin(2w)/w^2$, with values $1,0$ at zero, and $X=Kh$.
The modulation is
\[
\Sigma_2=\frac2{\sin\theta}
[\sin(\theta d)\Sigma_c-\cos(\theta d)\Sigma_s],\quad
\Sigma_\bullet=\frac{2\varepsilon}{\pi}
\left[\mathcal T_\bullet+\frac h2g_\bullet(X)\right],
\]
where $\mathcal T_\bullet$ is the trapezoidal rule on $[0,X]$
with mesh $h$. Its half-weight at zero is the exact diagonal contribution;
it cannot be discarded. Twice integrating the interpolation error on
each cell gives
$|\mathcal T_\bullet-\int_0^Xg_\bullet|
\le h^2\int_0^\infty|g_\bullet''|/8=O(h^2)$.
For finiteness, write $g_\bullet=\phi_\bullet/w^2$, where
$\phi_c=\frac12\cos2w-\frac14-\frac14\cos4w$ and
$\phi_s=\frac12\sin2w-\frac14\sin4w$:
the origin is removable and, for $w\ge1$,
$|g_\bullet''|\le6w^{-2}+8w^{-3}+6w^{-4}$.
Also $h|g_\bullet(X)|/2\le1/(2K)$.
Lemma~\ref{lem:two-integrals} gives
$\int_0^\infty g_c=0$, $\int_0^\infty g_s=\log2$ with tails at most $1/X$.
Consequently
$\Sigma_2=-2\pi^{-2}(\log2)\cos(\theta d)
 +O_\delta(\varepsilon^2+(N\varepsilon)^{-1})$,
using $\theta/\sin\theta=1+O_\delta(\varepsilon^2)$.
This proves \eqref{eq:hs-discrete}.
Finally $\Sigma_1=2\tr C_N(\varepsilon)^2$; divide the unmodulated
calculation by two and set $\varepsilon=2W$ to obtain
\eqref{eq:dpss-discrete}. That argument never divides by $\sin\theta$.
\end{proof}
The prolate constant is classical; \cite[eq.~(D.6)]{AbanovIvanovQian2011}
gives a sharper fixed-band expansion. At $\varepsilon=\tau/N$,
\eqref{eq:hs-discrete} matches the continuum logarithmic and oscillating
terms with error $O(\tau^2/N^2+1/\tau)$; this error does not vanish
at fixed $\tau$, whose convergence follows separately from the operator limit.

\subsection{Adjacent bands and cyclic words}\label{app:adjacent}
Put $C_\tau'=M_{e^{2\pi i\tau x}}C_\tau M_{e^{-2\pi i\tau x}}$,
$\mathcal N=M_{e^{-i\pi\tau x}}$ and $\mathcal V=\mathcal N^*C_\tau\mathcal N^*$.
With the positive-sign Fourier convention, the frequency bands of
$C_\tau,C_\tau'$ are $(-\tau/2,\tau/2)$ and $(-3\tau/2,-\tau/2)$.

\begin{proposition}[Adjacent products]\label{short:cor:adjacent-quadratic}
\label{T4:lem:factor}\label{rev5:prop:phase}\label{T4:conj:twoband}
Write $D(a)=\tr(C^{(a)}-(C^{(a)})^2)$,
$X=\tr(C_\tau^2C_\tau'^2)$,
$\Omega_2=\tr(\mathcal V^3\mathcal V^*)$, $\Omega_3=\tr\mathcal V^4$.
Then $\mathcal K_\tau=\mathcal V+\mathcal V^*$, the $\Omega_j$ are real, and
\begin{equation}\label{short:eq:cyclic}
\tr\mathcal K_\tau^4=2\tr C_\tau^4+4X+8\Omega_2+2\Omega_3.
\end{equation}
For every $\tau>0$,
\[
0\le X\le\tr(C_\tau C_\tau')=D(\pi\tau)-\tfrac12D(2\pi\tau)
=\frac{\log(2\pi\tau)+1+\gamma}{2\pi^2}+r,\quad
|r|\le\frac5{4\pi^3\tau},
\]
and $|\Omega_2|\le b(1/4)$, $|\Omega_3|\le4b(1/8)$.
Moreover $X=(12\pi^2)^{-1}\log\tau+O(1)$.
\end{proposition}
\begin{proof}
The kernel factorization is product-to-sum.
In degree four the cyclic word classes have sizes $2,4,8,2$:
alternating, contiguous two-and-two, three-and-one, and pure.
They give \eqref{short:eq:cyclic}; for example
$\tr(\mathcal V^2\mathcal V^{*2})=X$.
The antiunitary $f(x)\mapsto\overline{f(-x)}$ fixes $C_\tau,\mathcal N$
and $\mathcal V$, proving reality of the word traces.
The sum $C_\tau+C_\tau'$ is unitarily equivalent to $C_{2\tau}$,
so its defect gives the quadratic product identity.
Apply Lemma~\ref{lem:defect-exact} at $\pi\tau,2\pi\tau$ for the
error $1+1/4$ in units $(\pi^3\tau)^{-1}$.
The inequalities follow successively from $C_\tau^2\le C_\tau$
and $C_\tau'^2\le C_\tau'$ in positive trace pairings.

For $F(\phi)=\tr(e^{i\phi}\mathcal V+e^{-i\phi}\mathcal V^*)^4$,
the strip and trace comparisons give
$|F(\phi)-F(\psi)|\le16b(|\phi-\psi|/(2\pi))$.
The word expansion is
$F(\phi)=2\tr C_\tau^4+4X+8\Omega_2\cos2\phi+2\Omega_3\cos4\phi$.
Thus $16\Omega_2=F(0)-F(\pi/2)$ and
$8\Omega_3=F(0)+F(\pi/2)-2F(\pi/4)$, proving the bounds.
In degree two the same argument gives
\begin{equation}\label{T4:eq:osc1}
\Omega_1:=\tr\mathcal V^2=-\pi^{-2}(\log2)\cos(4\pi\tau)+\rho_1,
\qquad |\rho_1|\le3/(\pi^3\tau).
\end{equation}
Indeed $\tr\mathcal K_\tau^2=2\tr C_\tau^2+2\Omega_1$ and the
two quadratic errors sum to $6/(\pi^3\tau)$.

For the remaining coefficient, \cite[Theorem~2.1]{KitanineEtAl2009},
with $q=1,p(\lambda)=\lambda,F=1,g=0,x=2\pi\tau$, gives
$\log\det(I+\gamma C_\tau)
=2\tau\log(1+\gamma)+(2\pi^2)^{-1}\log^2(1+\gamma)\log\tau
+c(\gamma)+o(1)$.
The constant is analytic near zero; the fixed $n$th derivative of the
remainder there is $O((\log\tau)^n/\tau)$ by that paper's
eqs.~(8.40)--(8.43). The $\gamma^4$ coefficient $11/12$ therefore gives
$\tr C_\tau^4=2\tau-11(6\pi^2)^{-1}\log\tau+O(1)$, also consistent
with the Landau--Widom integral
$\int_0^1(t^4-t)/[t(1-t)]\,dt=-11/6$ \cite[p.~478]{LandauWidom1980}.
Subtract twice this formula from the quartic trace law for
$\mathcal K_\tau$ in \eqref{short:eq:cyclic} and divide by four.
\end{proof}
This computes one adjacent word, not arbitrary cyclic words, and does
not supply an independent proof of the quartic coefficient used.

